\documentclass[11pt,a4paper]{article}
\usepackage[margin=1.8cm,headheight=48pt,headsep=12pt,footskip=22pt]{geometry}
\usepackage{fontspec}
\setmainfont{Noto Serif}
\setsansfont{Noto Sans}
\setmonofont{DejaVu Sans Mono}
\usepackage{amsmath,amssymb,mathtools}
\usepackage{enumitem,tabularx,booktabs,array,multicol}
\usepackage{xcolor,graphicx,tikz,fancyhdr,lastpage}
\usepackage{listings,tcolorbox,hyperref}
\hypersetup{colorlinks=true,linkcolor=blue!60!black,urlcolor=blue!60!black}
\definecolor{hcmusblue}{RGB}{28,55,111}
\definecolor{lightblue}{RGB}{237,242,250}
\definecolor{codebg}{RGB}{247,248,250}
\setlength{\parindent}{0pt}
\setlength{\parskip}{4pt}
\setlist[enumerate]{leftmargin=*,itemsep=5pt,topsep=4pt}
\setlist[itemize]{leftmargin=*,itemsep=2pt,topsep=2pt}
\lstset{language=C++,basicstyle=\ttfamily\footnotesize,backgroundcolor=\color{codebg},frame=single,rulecolor=\color{black!15},breaklines=true,showstringspaces=false,columns=fullflexible,keywordstyle=\color{hcmusblue}\bfseries,commentstyle=\color{green!40!black},numbers=left,numberstyle=\tiny\color{black!45},xleftmargin=1.5em,framexleftmargin=1.2em}
\pagestyle{fancy}
\fancyhf{}
\fancyhead[L]{\includegraphics[height=1.05cm,keepaspectratio]{../../04_Nhan_dien/hcmus_logo.png}\hspace{0.45em}\begin{minipage}[c]{0.64\textwidth}\textbf{TRƯỜNG ĐẠI HỌC KHOA HỌC TỰ NHIÊN}\\[-1pt]\small ĐHQG TP. HỒ CHÍ MINH\end{minipage}}
\fancyhead[R]{\small\textcolor{hcmusblue}{TÀI LIỆU LUYỆN TẬP}}
\fancyfoot[L]{\scriptsize Biên soạn độc lập - Không phải đề thi chính thức}
\fancyfoot[R]{\scriptsize Trang \thepage}
\renewcommand{\headrulewidth}{0.5pt}
\renewcommand{\footrulewidth}{0.3pt}
\newcommand{\examtitle}[4]{\begin{center}{\Large\bfseries\color{hcmusblue}#1}\par\vspace{3pt}{\large\bfseries #2}\par\vspace{3pt}\small #3\quad\textbullet\quad #4\end{center}\vspace{5pt}\hrule\vspace{8pt}}
\newcommand{\studentline}{\begin{tabularx}{\textwidth}{X X}\textbf{Họ và tên:} \hrulefill & \textbf{MSSV:} \hrulefill\\[7pt]\textbf{Lớp:} \hrulefill & \textbf{Ngày làm:} \hrulefill\end{tabularx}\vspace{8pt}}
\newcommand{\points}[1]{\hfill\textbf{(#1 điểm)}}
\newtcolorbox{notebox}{colback=lightblue,colframe=hcmusblue!55,boxrule=0.5pt,arc=1.5mm,left=2mm,right=2mm,top=1mm,bottom=1mm}
\begin{document}
\examtitle{ĐỀ LUYỆN TẬP 01 - CƠ SỞ LẬP TRÌNH}{CƠ SỞ LẬP TRÌNH - HCMUS (theo lịch học tuần 1--5)}{90 phút - 10 điểm - C++11 - Không sử dụng tài liệu}{Giữa kỳ 1 - HK1 2026--2027}

\studentline
\begin{notebox}\textbf{Phạm vi:} Tổng quan C++, rẽ nhánh, vòng lặp, hàm, tham số mặc định/tham chiếu, mảng 1D và 2D. Dùng C++11; chỉ dùng các kiến thức đã nêu trong lịch học tuần 1--5.\end{notebox}
\begin{enumerate}
\item Xét chương trình sau. Hãy ghi chính xác kết quả in ra và giải thích thứ tự thay đổi của các biến. \points{2}
\begin{lstlisting}
int a = 2, b = 4, s = 0;
if (a++ < 3 && ++b == 5) s += b;
else s -= a;
for (int i = 1; i <= 3; ++i) s += i;
cout << a << " " << b << " " << s;
\end{lstlisting}
\item Viết hàm \texttt{sumMultiples(n, k)} trả về tổng các số nguyên dương không vượt quá $n$ và chia hết cho $k$. Đặt tham số mặc định $k=3$. Dùng vòng lặp và điều kiện phù hợp. Giả sử $n\ge0$, $k>0$. \points{2}
\item Cho đoạn mã:
\begin{lstlisting}
void adjust(int x, int &y) {
    x += 3;
    y += x;
}
int a = 2, b = 5;
adjust(a, b);
cout << a << " " << b;
\end{lstlisting}
Hãy cho biết kết quả in ra, giải thích sự khác nhau giữa truyền tham trị và truyền tham chiếu trong ví dụ. \points{2}
\item Viết hàm nhận mảng số nguyên $a$ có $n$ phần tử, trả về tổng các phần tử dương và đồng thời cập nhật biến tham chiếu \texttt{cnt} bằng số lượng phần tử dương. Nếu $n=0$, trả về 0 và \texttt{cnt}=0. \points{2}
Prototype: \texttt{long long sumPositive(const int a[], int n, int \&cnt);}
\item Viết hàm tìm chỉ số hàng có tổng phần tử lớn nhất trong ma trận $R\times C$. Nếu nhiều hàng cùng tổng lớn nhất, chọn hàng có chỉ số nhỏ nhất. Giả sử $R,C>0$. \points{2}
Prototype: \texttt{int rowWithMaxSum(const int a[][MAX], int R, int C);}
\end{enumerate}
\vfill\begin{center}\textit{--- Hết ---}\end{center}

\end{document}
